\section{2017: Attention Is All You Need}

\subsection{The Paper That Changed Everything}

In June 2017, a team of eight researchers at Google published a paper with a deliberately provocative title: \emph{``Attention Is All You Need''} \cite{vaswani2017attention}.

The paper introduced the \textbf{Transformer} architecture.
It dispensed entirely with recurrence --- no RNNs, no LSTMs, no processing one token at a time.
Instead, it proposed that \emph{attention} was sufficient to model language.

The results were immediate and dramatic: better performance on translation benchmarks, faster training, and a model that scaled gracefully with more data and more computation.

\subsection{What is Attention?}

Before Transformers, researchers had already been experimenting with attention as an add-on to RNNs \cite{bahdanau2014neural}.
The idea: when translating a word, instead of relying on a single compressed hidden state, let the model look back at all the input words and decide which ones are most relevant right now.

The Transformer took this idea and made it the entire foundation of the model.

\begin{tcolorbox}[colback=blue!5!white, colframe=blue!40!black, title=Attention in Plain English]
Imagine you are translating the French sentence ``Le chat noir dort'' into English.
When you are about to write the word ``black,'' your brain naturally focuses on ``noir'' --- not the whole sentence equally.

Attention does the same thing mathematically.
For each word being produced, the model computes a \emph{score} between that word and every word in the input, then uses those scores to create a weighted combination.
High-scoring words contribute more; low-scoring words contribute less.
\end{tcolorbox}

\subsection{Self-Attention: The Key Innovation}

The Transformer introduced \textbf{self-attention}: the model attends not just to an input sequence, but to its own sequence --- every word attends to every other word simultaneously.

This means the model can capture long-range dependencies in a single step.
The word ``it'' at the end of a long sentence can directly attend to the noun it refers to at the beginning, without the signal having to travel through dozens of sequential steps.

\begin{tcolorbox}[colback=green!5!white, colframe=green!40!black, title=Why Self-Attention Matters]
Consider: ``The trophy didn't fit in the suitcase because \textbf{it} was too big.''

What does ``it'' refer to? The trophy or the suitcase?
A human reader knows it's the trophy --- because if the suitcase were too big, it would have fit.

Self-attention lets the model directly compare ``it'' with every other word in the sentence and learn, from training data, that ``too big'' points back to the trophy.
An RNN has to carry this inference through sequential hidden states; self-attention resolves it in a single computation.
\end{tcolorbox}

\subsection{Parallelism: Training at Scale}

The other revolution was practical.
Because self-attention processes all positions simultaneously, Transformers are \emph{massively parallelizable}.
Every GPU core can work on a different part of the computation at the same time.

This meant that for the first time, language models could be trained on enormous datasets in reasonable time.
Scale was no longer the enemy --- it became the strategy.

Researchers quickly discovered an empirical pattern: \emph{bigger models trained on more data consistently performed better}.
This observation, later formalized as \textbf{scaling laws} \cite{kaplan2020scaling}, would drive the next phase of the field.
